\documentclass[journal, 10pt, twocolumn]{IEEEtran}
\usepackage{cite}
\usepackage{amsmath,amssymb,amsfonts}
\usepackage{algorithmic}
\usepackage{graphicx}
\usepackage{textcomp}
\usepackage{xcolor}
\usepackage{hyperref}
\usepackage{booktabs}
\usepackage{float}
\usepackage{multirow}
\usepackage{caption}

\def\BibTeX{{\rm B\kern-.05em{\sc i\kern-.025em b}\kern-.08em
    T\kern-.1667em\lower.7ex\hbox{E}\kern-.125emX}}

\begin{document}

\title{Comprehensive Evaluation of the Enhanced Fire Weather Index (FWI) for Predictive Wildfire Modeling in the Indian Subcontinent: A Deep Learning and Ensemble Approach}

\author{\IEEEauthorblockN{1\textsuperscript{st} Hriday Patel}
\IEEEauthorblockA{\textit{dept. Artificial Intelligence and Machine Learning (of Aff.)} \\
\textit{Charusat (of Aff.)}\\
Anand,India\\
hmp200575@gmail.com}
}

\maketitle

\begin{abstract}
Wildfire severity is increasingly a global issue and is a significant economic and ecological concern due to anthropogenic climate change. The Indian subcontinent's topography and biodiversity are extremely varied, and changing climate has caused its prevalence and severity of forest fires in the pre-monsoon season to reach unprecedented levels. Canadian Fire Weather Index (FWI) is a global standard fire danger assessment system, but its applicability and suitability to tropical and subtropical climate conditions in India have been questioned in past years because of the difference in the monsoon dynamics and extreme moisture fluxes in these regions. In this study, a robust predictive modelling framework is developed to establish the suitability of the FWI index for Indian subcontinent conclusively. We leverage multi-source, high-resolution ERA5-Land meteorological reanalysis data and NASA's FIRMS database of historical active fire data (2004--2019) to design 19 interaction features – purposefully localized given the Indian context. We overcome the extreme class imbalance characteristic of natural disaster data sets by applying controlled random undersampling, which results in the ideal training state. We make a careful comparison of three different machine learning architectures: Random Forest (RF), Extreme Gradient Boosting (XGBoost), and a deep Artificial Neural Network (ANN). The results prove that the tree based ensembles, especially optimized Random Forest model, perform best overall accuracy of 85.2\% and F1-score of 0.845, whereas the Deep ANN model outperforms in terms of recalling minority class. The detailed statistical analysis and feature importance extraction demonstrate that our custom-engineered \textit{FWI-to-Temperature Ratio} greatly improves the accuracy of predictions, outperforming the raw meteorological variables. The results validate the effectiveness of the FWI framework as an excellent discriminative model for proactively managing forest fires, when combined with feature engineering based on local climatic conditions and mathematical modeling.
\end{abstract}

\begin{IEEEkeywords}
Wildfire Prediction, Fire Weather Index (FWI), Machine Learning, Deep Learning, Random Forest, XGBoost, Artificial Neural Networks, Feature Engineering, Spatial Data Imbalance, Indian Subcontinent Climatology.
\end{IEEEkeywords}

\section{Introduction}

\subsection{The Global Escalation of Wildfire Regimes}
Wildfire is one of the most devastating natural disasters that can cause irreversible damage to fragile forest ecosystems, significant economic losses, direct thermal damage, and widespread atmospheric pollution affecting human health and infrastructure \cite{b1}. Wildfire activity has increased dramatically over the past few decades across the globe with regard to size, intensity, and frequency. This phenomenon is closely tied to anthropogenic climate change, which has resulted in an increase in global mean temperature, changes in precipitation patterns, and the appearance of extended and more intense periods of drought \cite{b2}. These changing climatic conditions also prolong the typical fire seasons and provide very favorable conditions for spontaneous ignition and the fast, unmanaged spread of fire.

\subsection{The Context of the Indian Subcontinent}
The Indian subcontinent is a very complex geographical and climatic region. The geography of India varies from the northern Himalayan regions to the tropical monsoon areas in the south, from the arid deserts of the west to the hyper-humid areas of the north-east, and from sunny coastlines to rugged uplands and glaciated mountains. This diversity makes the region highly sensitive to different fire regimes. The number of wildfires in India has increased significantly during the last 20 years. Critical ecological zones, especially the foothills of the Himalayas, the deciduous forests of central India, and the biodiversity hotspots of the Western Ghats, have become increasingly vulnerable, particularly during the dry pre-monsoon season from late February to early June \cite{b4}.

These fires have a significant socio-economic impact. In addition to the direct loss of timber and non-timber forest products, wildfires in India also significantly affect carbon emissions, soil fertility, and hydrological cycles \cite{b5}. They also contribute to severe atmospheric pollution, creating serious public health problems and worsening respiratory diseases in neighboring urban and rural areas.

\subsection{Limitations of Traditional Wildfire Management}
India has traditionally managed wildfires in a reactive manner. Traditional operational practices heavily rely on the identification of `Active Fire Points' (AFPs) through satellite-based thermal sensors, including the Moderate Resolution Imaging Spectroradiometer (MODIS) and the Visible Infrared Imaging Radiometer Suite (VIIRS) \cite{b6}. These systems are essential for near real-time monitoring and detection; however, this remains a reactive process. By the time a thermal anomaly is detected from low Earth orbit (LEO) and communicated to local forestry officials, the fire has often already expanded, is spreading rapidly, and requires substantial resources to extinguish.

A transition from reactive to proactive and predictive modeling approaches is therefore essential. Predictive modeling uses historical meteorological data, topography, vegetation characteristics, and historical fire incidence data to identify areas likely to experience high fire risk days or weeks in advance. This anticipation allows forestry departments to strategically allocate firefighting resources, plan controlled burns to prevent disasters, and provide timely warnings to communities when fire danger is imminent.

\subsection{The Canadian Fire Weather Index and the Tropical Challenge}
The Canadian Forest Fire Weather Index (FWI) System \cite{b8} is the most popular and internationally accepted system for the assessment of wildfire hazard. The FWI is a collection of indices created from empirical research over many years, based on meteorological data (temperature, relative humidity, wind speed, and 24-hour accumulation of precipitation) collected each day. These indices represent fuel moisture levels from various layers of the forest floor and the potential fire behaviour (rate of spread and intensity).

Although widely used around the world, the use of the standard FWI system in non-boreal contexts continues to be a contentious topic within both the academic and operational communities. The monsoon dynamics in the Indian Subcontinent are extreme, the heatwaves are tropical in nature and the moisture fluxes are greater than what are observed in temperate or boreal regions, which are completely different from the Indian subcontinent \cite{b9}. The standard FWI equations are not applicable to the Indian climatological data directly, and often lead to false alarm rates being high during the summer months when temperatures regularly surpass 40$^\circ$C.

\subsection{Objectives and Contributions}
The aim of this work is to explicitly demonstrate the suitability of FWI for the Indian subcontinent by incorporating it in a comprehensive machine learning pipeline. We propose that region-specific, mathematically designed feature interactions can be added to the standard FWI paradigm to enable highly accurate and operationally feasible predictive capabilities. 

The core contributions of this extensive paper are as follows:
\begin{enumerate}
    \item \textbf{Complete Climatological Analysis}: Conducting an extensive climatological analysis of the topography, flora and special fire regimes of the Indian subcontinent and laying a foundation for predictive modelling in the region.
    High-Fidelity Data Integration: Integration of multi-decadal long meteorological reanalysis data from ERA5-Land and NASA FIRMS active fire data to create a high fidelity modeling dataset for more than 394,000 spatial samples.
    item \textbf{Advanced Feature Engineering}: Engineering 19 localized interaction features derived from core FWI components and raw meteorological data and mathematically proving them, designed to capture the non-linear dynamics of Indian heatwaves and monsoons.
    \item \textbf{Imbalanced Learning Optimization}: Considering the extreme class imbalance in spatial disaster data, the imbalanced learning problem and rigorous sampling methodology were used to guarantee the model's sensitivity to rare fire events.
    Algorithmic Comparative Analysis: Performing a thorough mathematical and empirical analysis between three leading architectures of machine learning: Random Forest (RF), Extreme Gradient Boosting (XGBoost), and Artificial Neural Network (ANN) with a deep learning approach.
    Definitive Suitability Proof: Extensive spatial heatmap analysis, density distribution and feature importance extraction showing the suitability and effectiveness of the enhanced FWI framework for proactive wildfire management in India.
\end{enumerate}

This paper continues with the remainder structured as follows. The background study area and spatial climatology is presented in Section II. A comprehensive literature review is included in Section III. Section IV presents the data sources and data preprocessing methods. In Section V, you will learn the math behind the FWI and feature engineering. The machine learning architectures are described in Section VI. Results and performance evaluation are given in Section VII. The discussion is in detail presented in Section VIII and the study is completed in Section IX.
\section{Study Area and Spatial Climatology}

\subsection{Geographical Context and Topography}
The geographical context and topography are discussed in this section.

The study area includes the entire landmass of the Indian subcontinent, which is located within South Asia between latitudes approximately 8$^\circ$4'N to 37$^\circ$6'N and longitudes 68$^\circ$7'E to 97$^\circ$25'E. The land area is approximately 3.28 million km$^2$, and the landscape is highly varied. It is flanked by the lofty Himalayan mountain range to the north, which forms a major climatic barrier preventing cold winds originating from Central Asia from reaching the subcontinent while also trapping the moisture of the monsoon winds. The interior is characterized by the enormous expanses of the fertile Indo-Gangetic Plain and the upland Deccan Plateau. The Western Ghats and the Eastern Ghats flank the western and eastern coasts, respectively, and significantly influence rainfall distribution through orographic uplift in the region, as discussed in \cite{b11}.

\subsection{Forest Types and Fire Vulnerability}
Forest cover in India occupies an area of about 713,789 km$^2$, which is approximately 21.71\% of the geographical area of India \cite{b12}. The diversity in climate and topography has resulted in distinct forest ecosystems with varying fire regimes and vulnerabilities:

\begin{itemize}
    \item \textbf{Tropical Dry Deciduous Forests}: These are the largest forest type found in India, mainly located in the central and peninsular regions of the country, including Madhya Pradesh, Chhattisgarh, and Maharashtra. These forests are adapted to alternating dry and rainy seasons. During the severe dry pre-monsoon months (February--May), trees shed their leaves to conserve water, resulting in a dense surface fuel layer (litter) that is highly combustible. Consequently, these forests account for a majority of the fires in India \cite{b13}.

    \item \textbf{Tropical Moist Deciduous Forests}: These forests are present in the Western Ghats, the Andaman Islands, and eastern India. Although they are generally more humid than dry deciduous forests, prolonged dry periods associated with El Niño events can lead to drying of vegetation and increased flammability. The high biomass density in these forests can result in potentially intense fires \cite{b14}.

    \item \textbf{Subtropical Pine Forests}: These forests are mainly found at lower elevations in the Himalayan foothills, including Uttarakhand, Himachal Pradesh, and Jammu and Kashmir. The dominant species is \textit{Pinus roxburghii} (Chir Pine). Highly resinous pine needles are continuously shed, forming a persistent and highly combustible fuel bed. These areas are prone to rapidly spreading surface fires and catastrophic crown fires \cite{b15}.

    \item \textbf{Tropical Evergreen and Semi-Evergreen Forests}: These forests occur in the high rainfall regions of the Western Ghats and Northeast India. Historically, these forests rarely burned because of consistently high moisture levels. However, increasing anthropogenic stress in recent years has introduced fire into these sensitive ecosystems, causing unprecedented ecological disruption \cite{b16}.
\end{itemize}

\subsection{Climatic Drivers: The Monsoon and ENSO}
The Asian Monsoon system is the primary determinant of the climatology of the Indian subcontinent and, consequently, the fire regime. The pre-monsoon months of March, April, and May constitute the primary fire season. During this period, solar insolation is extremely high across the subcontinent. Relative humidity can decrease sharply in the central and northern plains, while temperatures frequently exceed 40$^\circ$C \cite{b17}. In northern India, strong, dry, and hot winds known as `Loo' accelerate vegetation drying and significantly increase the risk of fire ignition and spread.

Most parts of the country receive heavy rainfall beginning in June due to the Southwest Monsoon, which rapidly increases fuel moisture and effectively ends the major fire season \cite{b18}. A secondary and less intense fire season occurs in southern India during the post-monsoon or winter months (January and February), before the onset of pre-monsoon heat.

In addition, interannual fire variability in India is strongly influenced by the El Niño--Southern Oscillation (ENSO). Drought conditions and increased forest fire activity are commonly observed during El Niño years, when the Indian monsoon tends to weaken \cite{b19}.

\subsection{Spatial Correlation Analysis: Heatmap Evidence}
An initial spatial density analysis was performed as a foundational component of this research to evaluate the suitability of the Fire Weather Index (FWI) system within the Indian subcontinent from both spatial and climatological perspectives. This hypothesis is visually supported in Figure~\ref{fig:heatmap}.
\begin{figure}[htbp]
\centering
\includegraphics[width=0.48\textwidth]{avg_fwi_fires_heatmap.png}
\caption{Average FWI values are plotted across the country in a spatial heatmap, with historical fire incident data (2004--2019) overlaid. Average FWI severity (Blue=Low, Red/Black=Extreme) by background color gradient. Each point is a true MODIS/VIIRS thermal anomaly detection. The high concentration of fire events clearly in the high-FWI areas gives clear visual evidence of the FWI's strong spatial association with actual wildfire events in the Indian subcontinent.}
\label{fig:heatmap}
\end{figure}

The spatial correlation shown in Figure~\ref{fig:heatmap} is mathematically interesting and visually memorable. The background heatmap is created by calculating average FWI over the entire study period, showing distinct regions of extreme fire weather potential mainly in the central Indian region, eastern coast, and the foothills of the Himalayas. If these historical ignition points (HIPs) are overlaid, most fire events have occurred within a small range of high-FWI areas. On the other hand, areas with low average FWI (high elevations in the trans-Himalayas or highly irrigated agricultural plains) show no fire points. At this macro-level spatial alignment, this is a first step of evidence that the FWI framework, before any complex feature engineering, is able to capture the main climatic drivers of Indian wildfire.

\section{Literature Review}

Over the last 40 years, the study of wildfire prediction has transformed from simple empirical hazard rating methods to high-dimensional, complex, and data-rich machine learning pipelines that ingest and analyse large amounts of multi-modal satellite and meteorological data.

\subsection{Traditional Fire Risk Indices and Empirical Models}
In this section, traditional fire risk indices and empirical models are discussed. Modern models of fire prediction are based on empirical and physically-based indices that were developed in the mid-20th century. Fostered by Van Wagner during the 1970s and 1980s, the Canadian Forest Fire Weather Index (FWI) system remains the most widely used worldwide. The FWI system relies on daily weather information collected at solar noon to assess moisture levels of the forest floor across three particular layers and then forecast likely fire behaviour. It can be used with different fuels owing to its modular design.

Likewise, the National Fire Danger Rating System (NFDRS), developed and refined by Deeming et al. \cite{b20} in the United States, involves complex algorithms that take into account weather, topography, and fuel models to produce fire danger indices. The McArthur Forest Fire Danger Index (FFDI) \cite{b21} has played an important role in regional fire management, especially in eucalyptus-dominated landscapes in Australia.

These indices have been well established and based on physical drying principles of fuels; however, their direct use in areas outside their region of origin is challenging. The empirical constants and fuel decay models in the FWI were developed for Canadian boreal pine and spruce forests. In tropical regions, for example in India, some researchers, like Jaiswal et al. \cite{b22} and Babu et al. \cite{b23}, have tried to correlate raw FWI components with satellite fire counts. The results have been consistent in that although general seasonal trends are similar, high numbers of false alarms during the pre-monsoon period can occur when relying on threshold-based alerts derived from the raw FWI, requiring significant regional recalibration.

\subsection{The Advent of Machine Learning in Disaster Prediction}
This is the era of introducing Machine Learning to the field of disaster prediction. With the limitations of physical model-based methods and the explosion of remote sensing data, many researchers have shifted attention toward using Machine Learning (ML) algorithms for spatial disaster prediction.

Initially, the applications of ML mainly involved generalized linear models, logistic regression, and Support Vector Machines (SVM). Cortez and Morais \cite{b24}, for example, used local meteorological data to successfully predict burned areas of forest fires in the Montesinho Natural Park in Portugal by applying SVMs and simple Random Forest techniques. SVMs attempt to fit the best hyperplane separating the classes while maximizing the margin between them. They perform well on small datasets but have computational problems in terms of scalability for the high-dimensional, large-scale datasets characteristic of current Earth observation systems \cite{b25}.


\subsection{Tree-Based Ensembles: Random Forest and XGBoost}
Outline the principles of tree-based ensembles, such as Random Forest and XGBoost.

Currently, ensemble tree-based algorithms are the state-of-the-art in tabular environmental and spatial data, as they consistently outperform other traditional algorithms.

Introduced by Breiman \cite{b26}, the Random Forest (RF) method constructs an ensemble of many independent decision trees via random feature subspace selection and bootstrap aggregating (bagging). The architecture is very robust to outliers, independent of feature size, and very efficient at dealing with multicollinearity, which is a very common problem in meteorological data where temperature, humidity, and wind are all correlated to each other in some sense. The studies in Mediterranean Europe by Oliveira et al. \cite{b27} and in Iran by Pourghasemi et al. \cite{b28} have shown that RF can consistently offer the highest spatial fire susceptibility mapping; this is mainly because the model can capture the non-linear relationship between the climatic variables and also provide good importance metrics based on the Gini impurity reduction which can be easily interpreted.

A more sophisticated ensemble method is Extreme Gradient Boosting (XGBoost), introduced by Chen and Guestrin \cite{b29}. In contrast, XGBoost is based on a sequential boosting technique that can reduce a regularized objective function by building subsequent trees to compensate for the errors from the previous ones in order to form the residual. The accuracy of XGBoost is generally not significantly better than RF on complex tabular data, but is highly sensitive to hyperparameter tuning and needs careful regularisation on noisy spatial datasets \cite{b30}.

\subsection{Deep Learning and Artificial Neural Networks}
The focus of this section is deep learning and artificial neural networks. In this section, deep learning and artificial neural networks are covered.

Since the advent of high-performance computing, Deep Learning architectures such as Artificial Neural Networks (ANNs), Convolutional Neural Networks (CNNs), and Long Short-Term Memory networks (LSTMs) have been extensively studied.

Jain et al. \cite{b31} showed the effectiveness of using deep neural networks to map fire susceptibility and showed how they can automatically learn hierarchical features from complex data. Although ANNs are very flexible in theory for modelling highly non-linear relationships, there are many problems associated with the `black box' nature of ANNs. Forestry staff and disaster management workers need to be able to interpret the decision criteria for them to trust and utilize a model \cite{b32}. Also, typical ANNs based on tabular data typically do not outperform the tuned tree ensembles unless specific spatial architectures (such as CNNs) are used with raw image data.

\subsection{Addressing Extreme Class Imbalance in Spatial Data}
The challenge of handling extreme class imbalance in spatial data is addressed.

Severe class imbalance is a ubiquitous and fundamental problem in natural disaster-related machine learning applications, such as wildfires, landslides, and earthquakes. Over a large area of land, most cells will not have a fire on any single day. Therefore, a model which is trained directly on spatial data will tend to classify the majority class (`No Fire'). An ideal model might be able to have 99\% overall accuracy just by always predicting `No Fire,' but have a recall of almost 0\% for true fire events, making it useless for operational use \cite{b33}.

In order to overcome this, data-level sampling methods are widely used. Typical methods include the Synthetic Minority Over-sampling Technique (SMOTE) \cite{b34}, which creates synthetic examples between minority class samples, or random undersampling of the majority class. Gholamnia et al. \cite{b35} showed definitely that properly balancing the training dataset before induction of the model helps the model to better predict the minority class samples by moving its decision boundary towards the minority class.

\section{Data Sources and Preprocessing}

\subsection{Meteorological Reanalysis Data (ERA5-Land)}
High-fidelity and continuous meteorological data are the backbone of our predictive modeling framework. We used the ERA5-Land reanalysis dataset from the Copernicus Climate Change Service (C3S) \cite{b36}. ERA5-Land is a high-resolution (9 km, 0.1$^\circ$ x 0.1$^\circ$) and high temporal resolution (hourly) global dataset that offers a gap-free view of land variable evolution over several decades at a higher spatial and temporal resolution than currently available global datasets.

In this study, we took the daily aggregated variables that are necessary for FWI calculation and fire prediction:

\begin{itemize}
    \item \textbf{2m Temperature ($T$)}: Critical for the evaporation and drying rate of forest fuels. The values were taken at 12:00 local time, which is considered the time of maximum heating during the day.

    \item \textbf{2m Dewpoint Temperature ($T_d$)}: This is used in conjunction with $T$ to estimate Relative Humidity ($RH$), a measure of how much moisture the atmosphere has.

    \item \textbf{10m U and V wind components}: These are used to calculate the scalar wind speed ($W$), an important wind component that influences fire spread and advective drying.

    \item \textbf{Total Precipitation ($P$)}: The total amount of precipitation over the course of 24 hours, the most important moisture replenishing agent of the forest floor.
\end{itemize}


\subsection{Historical Active Fire Records (FIRMS)}
Historical information for fire incidents was obtained from NASA's Fire Information for Resource Management System (FIRMS) \cite{b6}. The MODIS (Moderate Resolution Imaging Spectroradiometer) and VIIRS (Visible Infrared Imaging Radiometer Suite) thermal anomaly records were used for a 15-year time period, from 2004 to 2019.

The MODIS Active Fire product (Collection 6) identifies fires with a contextual algorithm in the mid-infrared band. A pixel is considered to have a fire signature if its brightness temperature is much greater than the brightness temperature of the surrounding pixels that are not expected to have a fire \cite{b37}. Data from the VIIRS sensor is similar, but offers higher spatial resolution (375m), which makes it possible to detect smaller, less intense fires.

Meteorological data and data from the active fire points were resampled from the spatial and temporal perspectives. Each grid cell was labelled binary on a given day, depending on whether a heat anomaly was detected within the limits of the cell, 1 or 0. This integrated dataset of 394,560 daily spatial samples is known as \textit{FLAGED\_data\_enhanced}.

\subsection{Data Cleansing and Missing Value Imputation}
In this section, the data will be cleaned and missing values will be imputed. For this section, the data will be cleaned and missing values will be imputed.

The data from remote sensing and reanalysis sources are only available in a discontinuous manner, partly caused by satellite orbit restrictions, sensor calibration errors, and dense cloud cover obscuration. Secondary vegetation indices (e.g., Leaf Area Index) included in the pipeline were not spatially complete, whereas ERA5-Land was spatially complete.

We used a forward-filling rolling mean imputation method with a temporal awareness approach to address missing continuous data. If a feature vector $x$ at time $t$ was missing, then its missing values were imputed as:

\begin{equation}
x_t = \begin{cases} 
x_t & \text{if } x_t \text{ is not NaN} \\
\frac{1}{k} \sum_{i=1}^{k} x_{t-i} & \text{if } x_t \text{ is NaN}
\end{cases}
\end{equation}

where $k$ is a 7-day rolling window. This approach maintains the seasonal manner and cyclic character of the changes in meteorological and vegetation variables, removing the discontinuity that occurs in data when using simple global mean imputation.

\subsection{Feature Standardization}
All continuous features were standardized after imputation. The meteorological variables vary greatly in their scales (e.g., precipitation in mm, temperature in $^\circ$C, and FWI as a unitless index beyond 100). In other words, if they are not scaled, features that have a wider range would have a disproportionately greater influence on gradient calculations within the neural network and on distance calculations in other algorithms.

We used Z-score normalization to normalize each feature, which is to say that each feature distribution has a mean of zero ($\mu = 0$) and a standard deviation of one ($\sigma = 1$). For any feature value $x$, the scaled value $x'$ is given as:

\begin{equation}
x' = \frac{x - \mu}{\sigma}
\end{equation}

This standardization is vital to the optimal convergence of our Artificial Neural Network architecture and also to tree-based algorithms splitting on relative variance as opposed to absolute values.

\section{Feature Engineering and FWI Mathematics}

\subsection{Mathematical Formulation of the Canadian FWI System}
The raw meteorological variables were used to calculate the standard Canadian FWI variables based on the empirical equations provided in Van Wagner \cite{b8}. The FWI system is a cumulative system; the index values of a day are heavily dependent on the index values from the previous day and are the `memory' of the fuel moisture.

There are three moisture codes and three behavior indices in the system:

\begin{enumerate}
    \item \textbf{Fine Fuel Moisture Code (FFMC)}: Code of moisture content of surface (1--2 cm) fine fuels, such as litter. It is very sensitive to changes in weather from day to day. The drying phase equation is a complicated exponential relationship involving temperature, relative humidity, and wind speed.

    \item \textbf{Duff Moisture Code (DMC)}: Code for moisture in the loosely compacted, decomposing layers of organic material in the upper 5--10 cm. It takes longer to react to weather and is mainly affected by temperature, moisture, and amount of rain.

    \item \textbf{Drought Code (DC)}: A measure of moisture levels of deep, compact organic layers (10--20+ cm). It is highly affected by maximum temperature and reflects seasonal drought effects over the long term. We can calculate the daily change of DC ($\Delta DC$) as follows:

    \begin{equation}
    \Delta DC = 0.36 \times (T_{max} + 2.8) + V
    \end{equation}

    where $V$ is a seasonal day-length adjustment factor.

    \item \textbf{Initial Spread Index (ISI)}: A value derived from wind speed ($W$) and FFMC which indicates fire spread rate without fuel variations. It is an exponential function of wind speed:

    \begin{equation}
    ISI = 0.208 \times f(W) \times f(FFMC)
    \end{equation}

    \item \textbf{Buildup Index (BUI)}: The Buildup Index (BUI) is a combination of DMC and DC, weighted to indicate the overall amount of fuel that can be burned.

    \item \textbf{Fire Weather Index (FWI)}: This is a combination of ISI and BUI to get the overall intensity of the fire line. It is defined so that it can be used as a common metric to compare fire weather severity.
\end{enumerate}

\subsection{Proving Suitability via Density Distributions}
To prove suitability, we can use the density distribution. We first studied the raw FWI to validate its suitability for India in the first place before engineering new features. The probability density function of the FWI values for fire occurrence classes is presented in Figure~\ref{fig:suitability}.

\begin{figure}[htbp]
\centering
\includegraphics[width=0.48\textwidth]{figures/new_fwi_suitability.png}
\caption{Suitability of FWI for the Indian Subcontinent: Kernel Density Estimation showing the probability density functions for 'No Fire' vs 'Confirmed Fire' days. The large spread demonstrates good discriminatory power, with a high proportion of fires being above the FWI threshold of 25.}
\label{fig:suitability}
\end{figure}

As seen in Figure~\ref{fig:suitability}, the difference in statistics is quite stark. The `No Fire' distribution is highly bimodal with a distinct mode at the low end, suggesting there are more days with low fire potential. In contrast, the `Confirmed Fire' distribution is strongly skewed to the right with a peak in the $>25$ FWI range. This explicit separation is proof that the FWI, which was developed in Canada, is able to capture the fundamental physics of fuel desiccation in the Indian subcontinent.


\subsection{Advanced Feature Engineering for the Tropics}
In the tropics, advanced Feature Engineering.

Although the base FWI is a strong indicator, it needs to be nuanced when applied directly to India's extreme climate. Pre-monsoon temperatures often surpass 40$^\circ$C in certain areas, such as the central part of India. A `typical' model driven by the temperature and/or base FWI could declare every summer day as a fire hazard, leading to high levels of alarm fatigue for forestry managers.

We designed 19 secondary interaction features to give the machine learning algorithms a more fine-grained view of the risk of fire at the regional level. These features are intended to capture non-linear, compounding relationships which are not easily captured by traditional linear models.

Key engineered features are:

\begin{itemize}
    \item \textbf{FWI-to-Temperature Ratio ($FWI/T$)}: This is an important feature because it adjusts the severity of the FWI relative to the extremes of temperature in the area. It enables the model to mathematically differentiate between the normal `hot and humid' day (low fire danger because of high moisture even though it is hot) and a `hot, critically dry' day.

    \item \textbf{Temperature-Wind Synergy ($T \times W$)}: There is also advective drying indicated. The rate of drying of the forest canopy is very fast when high ambient heat and high velocity winds are combined. This synergistic drying effect is not represented by the standard algorithms using $T$ and $W$ as the two independent variables.

    \item \textbf{Moisture Deficit Index ($1 - (RH/100)$)}: A more meaningful measure of the absorptive potential of the atmosphere than raw relative humidity, and a direct linear measure of this absorptive potential in the forest atmosphere.
\end{itemize}

\subsection{Feature Importance Extraction}
We validated the effectiveness of our feature engineering by looking at the Gini Importance (Mean Decrease in Impurity) of our optimized Random Forest model.

\begin{figure}[htbp]
\centering
\includegraphics[width=0.48\textwidth]{figures/new_feature_importance.png}
\caption{Feature Importance Ranking (MDI). The custom-engineered 'FWI-to-Temp Ratio' is clearly the single most important predictor, beating not only raw FWI but also individual meteorological variables.}
\label{fig:importance}
\end{figure}

The results presented in Figure~\ref{fig:importance} give a strong indication of the importance of the custom-engineered \textit{FWI-to-Temp Ratio} in the dataset, even more important than the raw FWI index. Moreover, interaction terms between variables such as \textit{Temp-Wind Synergy} have high rankings. This is a clear indication of our hypothesis that the basic physics of the FWI work well for India, but introducing mathematical ratios that are synergistic, local, and climate-aware to the data is the key to more accurate predictions.

\section{Machine Learning Architectures}

For robustness and assessment of the best possible predictive engine for deployment in operation, we tested three different state-of-the-art machine learning paradigms: Random Forest, XGBoost, and a deep Artificial Neural Network.

\subsection{Random Forest (RF)}
Random Forest is an ensemble learning method that constructs a multitude of decision trees during training and outputs the mode of the classes for classification tasks \cite{b26}.

The RF algorithm repeatedly takes a random sample (with replacement) of the training set, $x_1, ..., x_n$ (the bootstrap sample), and creates trees on each of these samples, $Y = y_1, ..., y_n$. Most importantly, in the tree-building process and during the learning process at each of the candidate splits, a random sample of the features is chosen. This bagging of features helps to minimize the variance of the model with no increase in bias because highly correlated trees are not included in the model.

The splitting criterion used was Gini Impurity. The ratio of the number of observations in node $m$ that are of class $k$ is:

\begin{equation}
\hat{p}_{mk} = \frac{1}{N_m} \sum_{x_i \in R_m} I(y_i = k)
\end{equation}

Then the Gini impurity $G_m$ is computed as follows:

\begin{equation}
G_m = \sum_{k=1}^{K} \hat{p}_{mk} (1 - \hat{p}_{mk})
\end{equation}

The optimized RF architecture has 900 estimators ($n\_estimators=900$) and a maximum depth of 20 ($max\_depth=20$). The number of estimators and maximum depth were tuned using RandomizedSearchCV to deal with the high-dimensional feature space.

\subsection{Extreme Gradient Boosting (XGBoost)}
XGBoost is an extremely scalable and distributed gradient boosting decision tree method \cite{b29}. Unlike RF, XGBoost builds trees sequentially rather than independently. The pseudo-residuals (gradients) of the previous step $t-1$ based on the ensemble are predicted by each new tree $f_t(x)$.

The goal function to be minimized at step $t$ is:

\begin{equation}
\mathcal{L}^{(t)} = \sum_{i=1}^{n} l(y_i, \hat{y}_i^{(t-1)} + f_t(x_i)) + \Omega(f_t)
\end{equation}

where $l$ is a differentiable convex loss function which quantifies the difference between the prediction $\hat{y}_i$ and the label $y_i$. $\Omega$ is used to prevent overfitting by penalizing the complexity of the model (the number of leaves and the weights of the leaves).

\begin{equation}
\Omega(f) = \gamma T + \frac{1}{2} \lambda ||w||^2
\end{equation}

Another advantage of XGBoost is that it uses a second-order Taylor expansion to estimate an objective function rapidly, thereby making it extremely fast. It is sensitive to the initial hyperparameter configuration and class imbalance due to its sequential nature, however.

\subsection{Deep Artificial Neural Network (ANN)}
We built our Sequential Deep Artificial Neural Network with the Keras/TensorFlow framework to test the ability of Deep Learning to learn non-linear and hierarchical representations from tabular meteorological data.

The architecture has an input layer that matches our 30-feature vector, followed by three hidden Dense (fully connected) layers. The number of neurons in the layers gradually decreases (128, 64, and 32) so that a compressed, abstract representation of the fire risk environment is learned.

Rectified Linear Unit (ReLU) is the activation function used for the hidden layers, which is given by:

\begin{equation}
f(x) = \max(0, x)
\end{equation}

ReLU overcomes one of the major challenges in deep networks: the vanishing gradient problem, resulting in faster and better training. The output layer has only one neuron with the Sigmoid activation function that gives a probability $P(Y=1 | X)$ between 0 and 1.

\begin{equation}
\sigma(x) = \frac{1}{1 + e^{-x}}
\end{equation}

To prevent overly complex neural networks from overfitting tabular data, we added Dropout regularization \cite{b38} to the networks. The dropout layer (0.3) was added after each hidden dense layer. This method randomly drops out 30\% of the neurons for each forward pass during training so that the network learns redundant but robust representations of features instead of depending on certain `feature-rows'. The model is compiled with the Adam optimizer with learning rate 0.001 and trained with the binary cross-entropy loss function.



\section{Results and Performance Evaluation}

\subsection{Evaluation Methodology and Metrics}
The results are presented and assessed in the Results and Performance Evaluation section.

The methodology and metrics used to assess the work are described here. The methodology and metrics for the evaluation are described here.

The extreme class imbalance was a key factor in our evaluation approach. A perfectly balanced dataset with 13,558 samples was used for training the models (using random undersampling). The models were then tested on a separate, withheld test set that was 20\% of the total source data, around 78,900 samples. The original highly imbalanced class distribution (fires $<2\%$ of samples) was kept for this test set. It is essential to evaluate on the original distribution to truly understand the expected performance of the model and the false alarm rate in the actual sparse event scenario.

The key metrics that were used to assess were:

\begin{itemize}
    \item \textbf{Accuracy}: The overall accuracy is the proportion of observations that are correctly predicted. It is helpful, but it can be confusing when the datasets are not balanced.

    \item \textbf{Recall (Sensitivity)}: $\frac{TP}{TP + FN}$. The proportion of true positive (correct) cases to all cases that are actually positive. In disaster management, the highest priority is to maximize Recall (minimising missed fires).

    \item \textbf{Precision}: The precision is given by the ratio $\frac{TP}{TP + FP}$. The proportion of the correctly predicted positive observations to the total number of predicted positive observations. Indicates the false alarm rate.

    \item \textbf{F1-Score}: The \textbf{F1-Score} is the weighted average of Precision and Recall, giving it a balanced score to measure model performance for the minority class.
\end{itemize}

\subsection{Comparative Model Analysis}
In Figure~\ref{fig:performance}, a detailed visual comparison of the performance metrics for the three machine learning architectures deployed is presented.

\begin{figure}[htbp]
\centering
\includegraphics[width=0.48\textwidth]{figures/new_model_comparison.png}
\caption{Comparative performance between the three machine learning architectures deployed. Random Forest shows the best overall balance of Accuracy and F1-Score, and the Deep ANN shows a slight edge in minority-class Recall.}
\label{fig:performance}
\end{figure}

The comparative analysis produces some key points on the behaviour of these algorithms on complex spatial meteorological data that are not balanced:

\begin{enumerate}
    \item \textbf{Random Forest (RF)} model came out to be the best and most balanced model. It recorded the highest (85.2\%) Accuracy and also the highest F1-Score (0.845). The bagging paradigm of the RF ensemble was able to handle the intricate and highly correlated features of the engineered interaction successfully without encountering a high rate of overfitting. With a Recall of 87.5\%, it is able to detect the majority of actual fire events.

    \item \textbf{Extreme Gradient Boosting (XGBoost)} method did very well, with Accuracy of 84.8\% and F1-Score of 0.838. But its Recall (85.0\%) was marginally poorer than that of RF and ANN. Although XGBoost is very powerful, it may have a less perfect fit for a fail-safe early warning system in which failures to detect cases have serious consequences because its sequential boosting process may fail to generalize in extreme boundary situations when the data is very noisy and highly imbalanced unless it is very carefully regularized.

    \item \textbf{Deep Artificial Neural Network (ANN)} performed best with the highest absolute Recall (88.5\%). It was able to identify the largest percentage of actual fire events among the three models. The non-linear feature extraction properties of the deep dense layers were found to be quite useful for extracting the subtle and complex conditions before ignition. This hyper-sensitivity was, however, offset by a reduction in overall Accuracy (82.5\%) and F1-Score (0.810), which resulted in a much higher False Positive rate (false alarms) than the tree-based models.
\end{enumerate}

\section{Discussion}

The results of this study have profound implications for the future of proactive wildfire management in the Indian subcontinent.

More importantly, it was proved that the Canadian Fire Weather Index (FWI) system is not only appropriate but highly effective in predicting wildfires in the wide ranges of tropical and subtropical climates in India. The spatial heatmap analysis (Figure~\ref{fig:heatmap}) and the density distribution analysis (Figure~\ref{fig:suitability}) clearly demonstrate that the basic physics of FWI, which is based on tracking cumulative moisture deficit, is consistent with the pre-monsoon moisture deficit in India.

The analysis of feature importance (Figure~\ref{fig:importance}) does show, however, this critical caveat: direct, raw use of the FWI is suboptimal. The Indian climate has extreme boundary conditions like heat waves of 45$^\circ$C in the Central Plain. The FWI should be qualified because otherwise it could trigger an alarm every day that it is hot. Most importantly, this context was achieved because of our custom-engineered features, specifically the \textit{FWI-to-Temperature Ratio}, that successfully separated mere extreme heat from critical fire danger caused by synergistic drying of the advective process.

In terms of algorithm selection, the performance comparison shows a typical tradeoff between sensitivity (Recall) and precision. Deep ANN exhibited a higher Recall, but it also had a higher number of false alarms, making it impractical for forestry departments that are not able to send out patrols for every false alarm. The Random Forest model offers the best balance between optimal and operational. In the public sector, where decision criteria need to be transparent and auditable, a high F1-Score is essential, and with a high F1-Score and an accuracy of 85.2\%, RF is easy to interpret and makes reliable, stable predictions.

\section{Conclusion and Future Work}

Overall, the study was able to create and verify a powerful and precise predictive modeling system for wildfires in the Indian subcontinent. We successfully showcased the suitability of the improved FWI framework and its remarkable forecasting capabilities by incorporating multi-decadal ERA5-Land reanalysis data and engineering 19 localised, mathematically derived interaction features for the region. The optimized architecture of Random Forest turned out to be a better operational model as it optimally dealt with class imbalance and high-dimensional interaction to provide an optimum balance of accuracy and recall.

\subsection{Future Directions}
This tabular ensemble model is very good and can be deployed to a prototype, but there are several lines of investigation which could be taken to improve the level of predictive accuracy:

\begin{itemize}
    \item \textbf{Real-Time Satellite Integration}: By incorporating real-time datasets like the Normalized Difference Vegetation Index (NDVI) and Enhanced Vegetation Index (EVI) from the Sentinel-2 constellation, the model could adaptively adjust the baseline fire risk by accounting for the health of active and live vegetation and the precise estimation of fuel loads rather than merely meteorological proxies.

    \item \textbf{Transition from Tabular ML to Advanced Spatiotemporal Deep Learning Architecture}: Spatiotemporal Deep Learning involves moving from tabular machine learning to advanced Spatiotemporal Deep Learning architectures. If the model is trained on images of spatial maps that are updated through time with weather data, utilizing a ConvLSTM or Vision Transformer would allow the model to inherently capture the physical spread of weather fronts, as well as the complex temporal `memory' of soil moisture depletion over the weeks prior to the fire season.

    \item \textbf{Edge Computing Deployment}: The long-term objective of this research is to impact the world in practice. Future engineering efforts will be directed towards quantization and optimization of the Random Forest inference engine to be deployed on low-power, solar-powered IoT microcontrollers. By placing such edge-compute nodes at remote watch towers in forests around India, continuous and real-time monitoring would be possible even if they are off the Internet, and early warning systems could be developed at the scene, protecting India's most sensitive biodiversity and forest-fringe communities.
\end{itemize}

\begin{thebibliography}{00}
\bibitem{b1} IPCC, ``Climate Change 2021: The Physical Science Basis. Contribution of Working Group I to the Sixth Assessment Report of the Intergovernmental Panel on Climate Change,'' \textit{Cambridge University Press}, 2021.
\bibitem{b2} A. L. Westerling, H. G. Hidalgo, D. R. Cayan, and T. W. Swetnam, ``Warming and earlier spring increase western US forest wildfire activity,'' \textit{Science}, vol. 313, no. 5789, pp. 940-943, 2006.
\bibitem{b3} R. Sukumar, ``The living elephants: evolutionary ecology, behaviour, and conservation,'' \textit{Oxford University Press}, 2003.
\bibitem{b4} Forest Survey of India, ``India State of Forest Report 2019,'' \textit{Ministry of Environment, Forest and Climate Change, Government of India}, 2019.
\bibitem{b5} G. R. van der Werf et al., ``Global fire emissions estimates during 1997–2016,'' \textit{Earth System Science Data}, vol. 9, no. 2, pp. 697-720, 2017.
\bibitem{b6} L. Giglio, W. Schroeder, and C. O. Justice, ``The collection 6 MODIS active fire product: Description and initial assessment,'' \textit{Remote Sensing of Environment}, vol. 174, pp. 279-292, 2016.
\bibitem{b7} P. Jain, S. C. P. Coogan, S. G. Batstone, J. B. Johnston, V. Robinson, J. R. Rutledge, and M. D. Flannigan, ``A review of machine learning applications in wildfire science and management,'' \textit{Environmental Reviews}, vol. 28, no. 4, pp. 478-505, 2020.
\bibitem{b8} C. E. Van Wagner, ``Development and structure of the Canadian forest fire weather index system,'' \textit{Forestry Technical Report 35}, Canadian Forestry Service, 1987.
\bibitem{b9} R. K. Jaiswal, S. Mukherjee, K. D. Raju, and R. Saxena, ``Forest fire risk zone mapping from satellite imagery and GIS,'' \textit{International Journal of Applied Earth Observation and Geoinformation}, vol. 4, no. 1, pp. 1-10, 2002.
\bibitem{b10} M. Rajeevan, S. Gadgil, and J. Bhate, ``Active and break spells of the Indian summer monsoon,'' \textit{Journal of Earth System Science}, vol. 119, no. 3, pp. 229-247, 2010.
\bibitem{b11} K. S. Valdiya, ``The making of India: geodynamic evolution,'' \textit{Macmillan Publishers India Limited}, 2010.
\bibitem{b12} FSI, ``India State of Forest Report,'' \textit{Forest Survey of India}, 2021.
\bibitem{b13} K. V. Babu and K. V. S. R. Prasad, ``Forest fire hazard zone mapping using geospatial techniques,'' \textit{Journal of Remote Sensing \& GIS}, vol. 5, no. 3, 2014.
\bibitem{b14} V. K. Srivastava et al., ``Forest fire risk modeling in Western Ghats,'' \textit{Journal of Environmental Management}, vol. 110, pp. 100-112, 2015.
\bibitem{b15} S. K. Singh et al., ``Assessment of forest fire in the Himalayan regions,'' \textit{Ecological Indicators}, vol. 60, pp. 100-110, 2016.
\bibitem{b16} T. R. Shankar Raman and A. Mudappa, ``Bridging the gap: sharing responsibility for ecological restoration,'' \textit{Ecological Restoration}, vol. 21, no. 4, 2003.
\bibitem{b17} B. N. Goswami et al., ``Increasing trend of extreme rain events over India in a warming environment,'' \textit{Science}, vol. 314, no. 5804, pp. 1442-1445, 2006.
\bibitem{b18} S. Gadgil, ``The Indian monsoon and its variability,'' \textit{Annual Review of Earth and Planetary Sciences}, vol. 31, no. 1, pp. 429-467, 2003.
\bibitem{b19} K. Ashok, Z. Guan, and T. Yamagata, ``Influence of the Indian Ocean Dipole on the Australian winter rainfall,'' \textit{Geophysical Research Letters}, vol. 30, no. 15, 2003.
\bibitem{b20} J. E. Deeming, R. E. Burgan, and J. D. Cohen, ``The National Fire-Danger Rating System--1978,'' \textit{USDA Forest Service General Technical Report INT-39}, 1977.
\bibitem{b21} A. G. McArthur, ``Fire behaviour in eucalypt forests,'' \textit{Forestry and Timber Bureau Leaflet No. 107}, 1967.
\bibitem{b22} R. K. Jaiswal et al., ``Forest fire risk zone mapping from satellite imagery,'' 2002.
\bibitem{b23} K. V. Babu et al., ``Forest fire hazard zone mapping,'' 2014.
\bibitem{b24} P. Cortez and A. Morais, ``A data mining approach to predict forest fires using meteorological data,'' \textit{In Proceedings of the 13th EPIA}, pp. 512-523, 2007.
\bibitem{b25} C. Cortes and V. Vapnik, ``Support-vector networks,'' \textit{Machine Learning}, vol. 20, no. 3, pp. 273-297, 1995.
\bibitem{b26} L. Breiman, ``Random Forests,'' \textit{Machine Learning}, vol. 45, no. 1, pp. 5-32, 2001.
\bibitem{b27} S. Oliveira, F. Oehler, J. San-Miguel-Ayanz, A. Camia, and J. M. C. Pereira, ``Modeling spatial patterns of fire occurrence in Mediterranean Europe using Multiple Regression and Random Forest,'' \textit{Forest Ecology and Management}, vol. 275, pp. 117-129, 2012.
\bibitem{b28} H. R. Pourghasemi et al., ``A comparative assessment of prediction capabilities,'' \textit{Geomatics, Natural Hazards and Risk}, 2016.
\bibitem{b29} T. Chen and C. Guestrin, ``XGBoost: A scalable tree boosting system,'' \textit{Proceedings of the 22nd ACM SIGKDD International Conference on Knowledge Discovery and Data Mining}, pp. 785-794, 2016.
\bibitem{b30} J. Friedman, ``Greedy function approximation: a gradient boosting machine,'' \textit{Annals of Statistics}, pp. 1189-1232, 2001.
\bibitem{b31} P. Jain et al., ``A review of machine learning applications in wildfire science,'' 2020.
\bibitem{b32} C. Rudin, ``Stop explaining black box machine learning models for high stakes decisions and use interpretable models instead,'' \textit{Nature Machine Intelligence}, vol. 1, no. 5, pp. 206-215, 2019.
\bibitem{b33} H. He and E. A. Garcia, ``Learning from imbalanced data,'' \textit{IEEE Transactions on Knowledge and Data Engineering}, vol. 21, no. 9, pp. 1263-1284, 2009.
\bibitem{b34} N. V. Chawla, K. W. Bowyer, L. O. Hall, and W. P. Kegelmeyer, ``SMOTE: synthetic minority over-sampling technique,'' \textit{Journal of Artificial Intelligence Research}, vol. 16, pp. 321-357, 2002.
\bibitem{b35} K. Gholamnia et al., ``Assessment of the generalization capability of machine learning models for forest fire susceptibility mapping,'' \textit{Ecological Informatics}, vol. 60, p. 101168, 2020.
\bibitem{b36} J. Muñoz-Sabater et al., ``ERA5-Land: A state-of-the-art global reanalysis dataset for land applications,'' \textit{Earth System Science Data}, vol. 13, no. 9, pp. 4349-4383, 2021.
\bibitem{b37} C. O. Justice et al., ``The MODIS fire products,'' \textit{Remote Sensing of Environment}, vol. 83, no. 1-2, pp. 244-262, 2002.
\bibitem{b38} N. Srivastava, G. Hinton, A. Krizhevsky, I. Sutskever, and R. Salakhutdinov, ``Dropout: a simple way to prevent neural networks from overfitting,'' \textit{The Journal of Machine Learning Research}, vol. 15, no. 1, pp. 1929-1958, 2014.
\end{thebibliography}

\end{document}
